\section{基礎}
さまざまな基礎や典型問題についてまとめる.

\subsection{生成される空間}

ベクトルがいくつか与えられて, それらで張られる空間の次元を求めるには, \tf{そのベクトルを行ベクトルに持つ行列を基本変形した行列のrankを計算するとよい.} rankは基本変形で不変な量である.

\begin{egprob} (2021基礎数学試験-2)
実数$a$に対し, ベクトル
\[v_1 = \bmat{a\\1\\1\\2}, v_2 = \bmat{2\\1\\0\\3}, v_3 = \bmat{3\\1\\-1\\4}, v_4 = \bmat{5\\1\\-3\\a+5}\]
で生成される$\mb{R}^4$の部分空間の次元を求めよ.
\end{egprob}
\begin{egans}
\bealn{
&A = \bmat{
2&1&0&3\\
3&1&-1&4\\
a&1&1&2\\
5&1&-3&a+5
}
\to
\bmat{
2&1&0&3\\
1&0&-1&1\\
a-2&0&1&-1\\
3&0&-3&a+2
}\\
&\xrightarrow{2行で掃}
\bmat{
2&1&0&3\\
1&0&-1&1\\
a-1&0&0&0\\
0&0&0&a-1
}
}
よって$a=1$のときrankは2で$a\neq 1$のときrankは4である.
\end{egans}
\begin{rem}
最短の基本変形を目指すのもよいが, いつもそうしていると時に計算を間違える. やはり行基本変形はやり方が決まっているぶん安定するというものなので, そのへんはバランスを取りたい.
\end{rem}




\subsection{さまざまな行列}
まず基本的な行列の名前について覚えておこう. 後にも取り上げる事項もあるが, 固有値や対角化可能性についてもまとめておく.
\ben
\item 直交行列とは${}^{t}A\cdot A = I_n$を満たす行列のこと. [例] \tf{回転行列} $\sumib{\cos{\theta} & -\sin{\theta} \\ \sin{\theta} & \cos{\theta}}$. この行列は実係数とされることが多い(?). \tf{固有値の絶対値は1.} 実成分であるとすれば, \tf{行列式}は$\pm 1$.  Euclid内積において\tf{各行は正規直交基底をなし, 各列も正規直交基底をなす.}
\item ユニタリ行列とは${}^{t}\ovl{A} \cdot A = I_n$を満たす行列のこと. \tf{固有値は絶対値1の複素数}. 行列式も絶対値1の複素数.  Hermite内積において\tf{各行は正規直交基底をなし, 各列も正規直交基底をなす.}
\item 対称行列とは${}^{t}A = A$を満たす行列のこと. Euclid内積で$\inpr{Au}{v} = \inpr{u}{Av}$とできる. \tf{直交行列で対角化可能.}(固有ベクトルの正規直交化を並べた行列と同値.)
\item Hermite行列とは${}^{t}\ovl{A} = A$を満たすものをいう. $\inpr{Au}{v} = \inpr{u}{Av}$ 左辺を$A^{\ast}$と書いたりすることもある. \tf{ユニタリ対角化可能.}(固有ベクトルの正規直交化を並べた行列と同値)
\item 正定値行列とは, $A\in M_n(\mb{R})$であって, 任意の$0$でない列ベクトル${\bf x}$に対して, ${}^{T}{\bf x}A{\bf x}$が必ず正となるとに言う.
\item 正規行列とは$A^{\ast}A = AA^{\ast}$を満たすもの. (しかし正直, どこで使われてるのかわからない)
\item
\een

\begin{rem}
いくつかの行列は\tf{内積と非常に相性が良い.}
\end{rem}

\begin{egprob}
$D\in M_{n}(\mb{R})$が歪Hermite行列, つまり$D^{\ast} = -D$を満たすとするとき, $D$の固有値は0か純虚数であることを示せ。とくに実歪対称行列$D$に関して主張が成り立つ($D^{\ast} = {}^{T} D$だから)。
\end{egprob}
\begin{egans}
$Dv = \lambda v$とする($v\neq \ve{0}$)。$\mb{C}^{n}$の標準内積に関して$\inpr{Dv}{v} = \inpr{v}{D^{\ast}v}$が成り立ち, $D^{\ast} = -D$だから
\beqn
&&\lambda\inpr{v}{v} = \inpr{\lambda v}{v} = \inpr{Dv}{v}\\
&=& \inpr{v}{D^{\ast}v} = -\inpr{v}{Dv} = -\inpr{v}{\lambda v}\\
&=& -\ovl{\lambda} \inpr{v}{v}
\eeqn
よって$\lambda = -\ovl{\lambda}$なので実部が0でなければならない。\qed
\end{egans}
\begin{rem}
元々, 歪Hermite行列だとか対称行列だとか直交行列だとか, いくつかの行列はその「行列の形」ではなく\tf{内積と絡めて}理解したほうがよい対象であったりする。たとえば, 実数上の内積では
\[\lan Av, w\ran = \lan v, {}^{t}Aw\ran\]
という関係式があったのだから, $\lan Av, Aw\ran = \lan v,w\ran$を満たすもの(つまり, 「直交性(内積0)を保つ行列」)として直交行列を定義すると思うと自然である。
\end{rem}
\begin{prac}
対称行列の固有値が実数であることを示せ。
\end{prac}
\begin{prac}
実の成分をもつ直交行列の固有値は, 絶対値1の複素数であることを示せ.
\end{prac}
\begin{warning}
固有値が$\pm 1$になるとは限らない! おそらく$\det{A} = \pm 1$との勘違い.
\end{warning}



\subsection{固有値}

\begin{prop} (トレースと行列式の性質)\\
\ben
\item $\tr{A+B} =\tr{B+A}$
\item $\tr{AB} = \tr{BA}$
\item $\det{AB} = \det{BA}$
\item $\tr{A}$は$A$の固有値和である。
\item $\det{A}$は$A$の固有値積である。
\een
\end{prop}



\begin{egprob} (Quiz, by Shakayami)
以下の行列の固有値を求めよ。
\[\bep
4 & 9 & 2\\
3 & 5 & 7\\
8 & 1 & 6
\enp
\]
%魔法陣ですね
%魔法陣なので、(1,1,1)が固有ベクトルになります(by shakayami)
\end{egprob}
\begin{egans}
行列を$A$として$A- \lambda I$を基本変形したとき
\[
\bep
15-\lambda & 15-\lambda & 15-\lambda\\
3 & 5-\lambda & 7\\
8 & 1 & 6-\lambda
\enp
\]
だから$\lambda = 15$は固有値\footnote{この行列は魔法陣であることに気づくと$(1,1,1)$が固有ベクトルとわかってそこから$15$が固有値であると気づくこともできる。}。$\lambda\neq 15$のとき上を$15-\lambda$で割って基本変形すると
\[
\bep
1 & 1 & 1\\
0 & 2 - \lambda & 4\\
0 & -7 & -2 - \lambda
\enp
\]
なので行列式を取ると$(2-\lambda)(-2-\lambda) + 28 = \lambda^2 + 24$なので$\lambda = \pm 2\sqrt{6}i$が固有値である。よって
\[\lambda = 15, \pm 2\sqrt{6}i\]
\end{egans}
上では15という非負実数の固有値を持っていた。実は次のようなことが知られている。線形代数の定理であるけれども, 実は幾何的な証明が知られている。
\begin{theo} (\tf{Frobeniusの定理})\\
各成分が負でない実数からなる$n$次正方行列は, 少なくとも一つ負でない実数の固有値をもつ。
\end{theo}










\subsection{次元}
\begin{theo}
$\dim{V} + \dim{W} = \dim{(V\cap W)} + \dim{(V+W)}$
\end{theo}
\begin{eg}
$n$次元ベクトル空間$V$の$n-1$次元ベクトル空間$W_1,W_2$に対して,$W_1\cap W_2$の次元を求める。$W_1\neq W_2$であるならば$W_1+W_2 = V$なので$W_1\cap W_2$の次元は定理より$n-2$である。$W_1=W_2$ならば明らかに$n-1$次元である。
\end{eg}
線型代数でも準同型定理は大切である.
\begin{prop}
線形写像$f:V\to W$があるとき$V/\ker{f} \cong \im{f}$. とくに, $\dim{V} - \dim{\ker{f}} = \dim{\im{f}}$である.
\end{prop}

理論的に使うことがあるので, 次を示しておく.
\begin{egprob} \label{egprob:rank_ineq}
$n\times n$行列$A,B$に対して, 次を示せ.
\[\rank{A} + \rank{B} \leq \rank{AB} + n\]
\end{egprob}
\begin{egans}
$A|_{\im{B}}$に関して準同型定理を考えると
$\rank{AB} = \rank{B} - \dim{(\ker{A} \cap \im{B})} $だから, 示すべき不等式は
\[\rank{A} + \dim{(\ker{A}\cap \im{B})} \leq n\]
である. これは確かに正しい. なぜなら,
\[n = \rank{A} + \dim{\ker{A}}\]
と比較できるから. \qed
\qed
\end{egans}
この不等式は後に対角化可能性を論じるときに用いられる.
\begin{egprob}
$n\times n$行列$A$の最小多項式が重根を持たないとき, $A$の固有ベクトルよりなる基底が存在することを証明せよ.
\end{egprob}

\begin{egans}
$A$の相異なる固有値を$\lambda_1,\cdots, \lambda_s$とし, $W_{\lambda_i}$を固有値$\lambda_i$に関する固有空間とする. すなわち, $W_i = \ker{(A-\lambda_i I)}$である. このとき最小多項式は$\prod_{i=1}^{s}(X-\lambda_i)$となる. よって,
\[\prod_{i=1}^{s} (A-\lambda_i I) = O\]
により,
\[\rank{\parena{\prod_{i=1}^{s} (A-\lambda_i I)} }= 0\]
である. 先の不等式を$s-1$回用いると,
\bealn{
0=\rank{\parena{\prod_{i=1}^{s} (A-\lambda_i I)}} &\geq -(s-1)n + \sum_{i=1}^{s} \rank{(A-\lambda_i s)}\\
&= -\ao{ \sum_{i=1}^{s} (n-\rank{(A-\lambda_i I)}) + n}\\
&= -\sum_{i=1}^{s} (\dim{W_{i}}) + n \quad (\because 準同型定理)
}
よって$\dim{\oplus_{i=1}^{s} W_i} = \sum_{i=1}^{s} \dim{W_i} \geq n$が得られる. 各$W_i$の基底を選び, それらの和集合を考えればそれは$n$本以上のベクトルを含み, かつ一次独立である. よって, それは実際には$n$本であり, $\mb{R}^{n} \cong W_1\oplus \cdots \oplus W_s$となる. ここから固有ベクトルよりなる基底が作れる.
\end{egans}

次に刺される読者は多いのではなかろうか.
\tbox{
\begin{egprob}
$k=\mb{R},\mb{C}$とする. 任意の$k$成分の行列$A$に対し, $A$の行ベクトルの一次独立な本数$(\mr{row. rank}A)$と$A$の列ベクトルの一次独立な本数$(\mr{column. rank}A)$は一致することを示せ.
\end{egprob}
}{}
\begin{egans}
$A$を$m\times n$行列とする. $A$は線形写像$k^n \to k^m$を定める. 線形写像$\Phi:k^m \to (k^m)^{\ast} = \mr{Hom}_{k}(k^m,k)$を次で定める.
\[\Phi(y) = (k^n\ni x\mapsto \lan Ax,y \ran \in k)\]
ただし$x=(x_i), y=(y_i)\in k^m$に対し$\lan x,y\ran = \sum_{i} x_iy_i$と定める. このとき, 次が成立することを示す.
\lefba{
\begin{claim} $V = \im{A}\subset k^m$とするとき,
\[\ker{\Phi} = V^{\bot} = \ker{{}^{t}A}\]
ただし$V^{\bot} = \set{x\in k^m\mid \forall v\in V,  \lan x, v\ran = 0}$である.
\end{claim}
}
\prf $k^n$の標準基底を$e_1,\dots, e_n$として
\bealn{
y\in \ker{\Phi} &\iff \forall x\in k^n, \lan Ax,y\ran = 0 \\
&\iff \forall v=Ax \in V, \lan v,y\ran = 0 \\
&\iff \forall x\in k^n, \lan x,{}^{t}Ay\ran = 0 \\
&\iff {}^{t}Ay = 0\qquad (\because x=e_iとすればわかる) \\
&\iff y\in \ker{{}^{t}A}
}
より成立する. (\tf{Claim証明終})\\
これにより
\bealn{
\dim{\ker{\Phi}} &= \dim{V^{\bot}} \\
&= m - \dim{V} = m - \mr{row.rank}A\\
&= \dim{\ker{{}^t A}} = m-\dim{\im{{}^t A}}\\
&= m - \mr{colmun.rank} A
}
により示される.
\end{egans}

{\small
\begin{rem} 上の主張は$k$を一般の体にしても実はよいのだが, $\dim{V^{\bot}} = m - \dim{V}$に少し問題がある(それ以外は準同型定理や線形写像の話なので, 正しいはずである). もともとこの次元等式を示すにあたって, [佐武]などの教科書では$W\cap W^{\bot} = \set{0}$を示すであろう. この部分は$\mb{R},\mb{C}$に非退化なノルムが入ることを用いているので, 一般の体$k$では適用できない. \cite{SLA} Theorem6.4(145p)から次の定理を引用する.
\begin{theo} $k$を体, $V,V'$を$k$ベクトル空間として,
$V\times V'\to k$を双線型写像とし, $W$はこの双線型写像に関して$V'$と直交する元全体からなる$V$の部分空間とし(\tf{左での核}という), $W'\subset V'$も同様(\tf{右での核})とする. $V'/W'$が有限次元であるとき, 誘導される写像$V'/W'\to (V/W)^{\lor}$は同型写像である.
\end{theo}
さて, 先の問題に戻って$\Phi$はもとより$k^n\times k^m \to k$, $(x,y)\mapsto \lan Ax,y\ran$という双線型写像から得られるものである. 左での核は$\ker{A}$であり, 右での核は$V^{\bot}$である. よって定理と, 有限次元の場合には双対を取っても次元が変わらないことから$n-\dim{\ker{A}} = m - \dim{V^{\bot}}$であり, 左辺は$\dim{\im{A}} = \dim{V}$だから$\dim{V} = m - \dim{V^{\bot}}$が得られる.
\end{rem}
}



\subsection{逆行列・余因子行列}

\begin{prop} (\tf{クラメルの公式})

\end{prop}
次の問題はCramerの公式を覚える良い切っ掛けになるであろう。
\begin{prob}(筑波大学大学院システム情報工学研究科経営・政策科学専攻・社会システム工学専攻)\\
平面$\mb{R}^2$上の3点$(x_i,y_i)$ ($i=1,2,3$)は一直線上の点ではないとする。$x_1,x_2,x_3$が互いに異なる実数であるとき, 次の問いに答えよ。
\ben
\item $\det{\bep x_1^2 & x_1 & 1\\ x_2^2 & x_2 & 1\\ x_3^2 & x_3 & 1 \enp} \neq 0$であることを示せ。
\item $\det{\bep x_1 & y_1 & 1\\ x_2 & y_2 & 1\\ x_3 & y_3 & 1 \enp} \neq 0$であることを示せ。
\item 3点$(x_i,y_i)$ ($i=1,2,3$)を通る放物線
\[y=ax^2+bx+c\quad (a,b,c\in \mb{R}, a\neq 0)\]
が存在することを示せ。
\een
\end{prob}

\newpage
\subsection{基本変形}
基本変形は慎重にやろう. 基本変形で保たれるものは次の通りである.
\begin{prop} 正則行列$P$と(正方とは限らない)行列$A$を考える.
\ben
\item 積$PA$が定義できるとき, $\ker{PA} = \ker{A}$である. したがって\tf{行基本変形は$A\bolm{x} = \bolm{0}$の解空間を変えない. }
\item 積$AP$が定義できるとき, $\im{AP} = \im{A}$である. したがって, \tf{列基本変形は$A$の像を変えない.}
\item 積$PA$が定義できれば$\rank{A} = \rank{PA}$であり, 積$AP$が定義できれば$\rank{A} = \rank{AP}$である. したがって, $\rank$は基本変形で不変な量である.
\een
\end{prop}






\begin{egprob} (H31基礎科目7) \label{prob:H31kiso7}
次の行列式を求めよ.
\[
\bmat{
0 & 1 & 2 & 3 & \cdots & n-1 \\
1 & 0 & 1 & 2 & \cdots & n-2 \\
2 & 1 & 0 & 1 & \cdots & n-3 \\
3 & 2 & 1 & 0 & \cdots & n-4 \\
\vdots & \vdots & \vdots & \vdots & \ddots & \vdots \\
n-1 & n-2 & n-3 & n-4 & \cdots & 0
}
\]
\end{egprob}

\begin{egans}
$n$行目から$n-1$行目を引く. $n-1$行目から$n-2$行目を引く. この順で, $i+1$行目から$i$行目を引き, 最後に$2$行目から1行目を引く. すると,
\[
\bmat{
0 & 1 & 2 & 3 & \cdots & \aka{\bolm{n-1}} \\
\aka{\bolm{1}} & -1 & -1 & -1 & \cdots & -1 \\
1 & \aka{\bolm{1}} & -1 & -1 & \cdots & -1 \\
1 & 1 & \aka{\bolm{1}} & -1 & \cdots & -1 \\
\vdots & \vdots & \vdots & \vdots & \ddots & \vdots \\
1 & 1 & 1 & 1 & \cdots & -1
}
\]
を得る. $j$ 列目$(2\leq j\leq n)$に$1$列目をたすと, 1行目と1列目を除いた部分で「1」は「2」に変わり, 「$-1$」は「0」に変わる. こうして得られた行列の行列式は, もとの行列式とは変わっておらず, 上の赤字の部分の数をかけていくことで計算できる. そこには$n-1$が一つと, 1が一つと, 2が$n-2$個あり, 上の赤字の順番で掛けていくときの符号は$(-1)^{n-1}$なので, 答えは\bolm{(-1)^{n-1}2^{n-2}(n-1)}.
\end{egans}







\subsection{($\heartsuit$)相似}
群論でいうところの共役関係が相似である。相似であることによって最小多項式は一致し, det や tr, その他様々な性質が同期する。
\besha
\begin{theo}
$A,B\in \mr{Mat}(n,\mb{C})$が相似であることと, $A,B$のJordan標準形が並び替えを除いて一致することは同値である。
\end{theo}
\ensha
\begin{rem}
相似であることを示すより相似でないことを示すほうが易しい。実際,Jordan標準形を見るまででも無く,  相似であることによって保存されるものを見て分かるケースが多いからである。たとえば$A,B$が相似であるなら, $A,B$の最初多項式や固有多項式は等しいということが容易に分かる。

なお, この逆の主張は成り立たない。たとえば, $E_{ij}\in \mr{Mat}(4,\mb{C})$を$i$行$j$列のみが1であるような行列とすると, \[E_{12} + E_{34},\quad E_{12}\]
は同じ最小多項式を持つが, 相似ではない(どちらも異なるJordan標準形になっているからである)。
\end{rem}













\clearpage
\subsection{演習問題}
\subsubsection{Level A}


\begin{prob} (H20基礎1) \label{prob:H20B1}
次の二つの行列は相似であるか。
\[
\begin{pmatrix}
1 & 1 & 0\\
0 & 1 & 0\\
0 & 0 & 1
\end{pmatrix},\quad
\begin{pmatrix}
1 & 0 & 0\\
0 & -1 & 0\\
0 & 0 & 1
\end{pmatrix}
\]
\end{prob}

\begin{prob} (H21基礎2) \label{prob:H21B2}
$a,b$を複素数とし,
\[
A=\begin{pmatrix}
a & 1 & 6 & 8\\
1 & a & -1 & 2\\
0 & 0 & b & 7 \\
0 & 0 & 0 & 2
\end{pmatrix}
,\quad
B=\begin{pmatrix}
9 & 0 & 0 & 0\\
-7 & 2 & 0 & 0\\
4 & 1 & 7 & 0 \\
-8 & 5 & 3 & 4
\end{pmatrix}
\]
\ben
\item $A,B$の固有値を求めよ。
\item $P\in GL(4,\mb{C})$が存在して$B=PAP^{-1}$となるような$a,b$を求めよ。
\een
\end{prob}

\begin{prob} (H9数学I-1)
$A^2 + E_2 = O$を満たす二次実正方行列$A$の例を一つ与えよ. また, $B^2 + E_3 = O$を満たす3次実正方行列$B$は存在しないことを示せ.
\end{prob}

%================================Level B===============
\subsubsection{Level B}
\begin{prob} (H24数学I, 1) \label{prob:H24-I-1}
$A,B\in \mr{Mat}_{n}(\mb{C})$, 多項式$f(X)\in \mb{C}[X]$とする. $Af(B) = B$が成り立っているとする.次を証明せよ.
\ben
\item $f(B)$が正則ならば$A$と$B$は可換.
\item $f(B)$が正則でなければ$f(0) = 0$.
\een
\end{prob}




\subsubsection{Level C}
\begin{prob} (H9数学I-3)
$A,B$を実$n$次正方行列とする. $P\in \mr{GL}_n(\mb{C})$があって$B=PAP^{-1}$とするとき, $Q\in \mr{GL}_{n}(\mb{R})$があって$B=QAQ^{-1}$となることを示せ.
\end{prob}


\subsubsection{Hint・Answer}

\begin{probans}
\chatgpthint{固有多項式と最小多項式を比較してください。}
\end{probans}
\begin{probans}
No. 固有多項式を見よ.
\end{probans}
\begin{probans}
\chatgpthint{固有値ごとの固有空間とJordan標準形を調べてください。}
\end{probans}
\begin{probans}
$A=\sumib{0&1\\ 1&0}$. $(\det{B})^2 = -1$より存在しない.
\end{probans}
%=====================-Level B===================
\begin{probans}
(1):一般に, $AC = B$かつ$BC=CB$ならば$ABC=BAC$である. $C=f(B)$とせよ. \\
(2): $f(B)v = 0$なるノンゼロなベクトル$v$を取る.
\end{probans}

%=====================Level C====================
\begin{probans}
$R(z) := zP + \ovl{z}\ovl{P}$を考え, $\parena{\dfrac{d}{dz}}^n \det{R(z)} \neq 0$を示せ.
\end{probans}









%==================================================
\newpage
